% Integrated FF.41; review copy only, do not input separately.
\begin{claim}{FF.41}{Proved}\label{clm:FFmetabelianbarrier}
The invariant-lift step of
Claim~\ref{clm:FFmaxabelianreconstruct}
fails for an actual elliptic curve over
the metabelian field \(B_p\), even
after every finite change of local
arithmetic model. More precisely, put
\(H=I''=G_{B_p}\), and take
\[
                  A:\quad y^2+xy=x^3-p.
                                                        \tag{FF.41.1}
\]
There is a finite local field
\(M_0/\mathbb Q_p\) of degree at most
two, whose chosen global quadratic
model lies in \(K_p^{\rm ab}\), such
that the following holds. The space
\[
 S=\operatorname{Hom}_{H}
          (V_pA,\mathbb Q_p(1))
                                                        \tag{FF.41.2}
\]
is two-dimensional. Let \(N\) be the
subspace of functionals whose pushforward
of the point Kummer class is zero for
every \(a\in A(E)\) and every finite
extension \(E/B_p\). Then for every
finite extension \(M/M_0\),
\[
       N\simeq\mathbb Q_p(1),\qquad
       S/N\simeq\mathbb Q_p,\qquad
       S^{I_M}=0,\qquad
       (S/N)^{I_M}=S/N\ne0.
                                                        \tag{FF.41.3}
\]
The first two isomorphisms describe
the actual local Galois representations.
Consequently the visible invariant
class has no inertia-invariant lift.
This is a counterexample to a proposed
replacement lemma, not a counterexample
to \(\mathcal R(B_p)\).
\end{claim}
\begin{proof}
For the displayed integral Weierstrass
equation one computes
\[
             c_4=1,\qquad
             \Delta=p(1-432p),\qquad
             v_p(j)=-1.
                                                        \tag{FF.41.4}
\]
In particular \(j\ne0,1728\).
Tate's convergent inverse to the
\(j\)-series gives \(q\in\mathbb Q_p\)
with \(v_p(q)=1\) and a Tate curve
\(E_q\) with this \(j\)-invariant.
Two characteristic-zero elliptic curves
with this \(j\)-invariant become
isomorphic over a quadratic extension:
their geometric automorphism group
is \(\{1,-1\}\), so the isomorphism
torsor is quadratic. Thus
\(A_{M_0}\simeq E_q\) for a field
\(M_0\) of degree at most two.
A square class in \(\mathbb Q_p^\times\)
has a rational representative, by
openness of squares and density of
\(\mathbb Q\). Hence \(M_0\) is
the completion of a chosen quadratic
number field \(D_0\). Since its
inertia character is abelian,
\(D_0\subset K_p^{\rm ab}\).
One could also use the split node
with tangent cone \(y(y+x)\) in
(FF.41.1); allowing the quadratic
extension avoids needing a separate
split-reduction criterion.

The exact primary uniformization
input is Tate's Theorem~1, together
with the description of finite-order
points on pages~18--19 of the author
text \citep{FF_Tate1995}. It gives
the compatible local Galois sequences
\[
 \begin{gathered}
   \varphi:M^\times/q^{\mathbb Z}
                         \xrightarrow{\sim}A(M),\qquad
   0\longrightarrow\mathbb Z_p(1)
    \longrightarrow T_pA
    \longrightarrow\mathbb Z_p\longrightarrow0,\\
   0\longrightarrow\mathbb Q_p(1)
    \longrightarrow V_pA
    \longrightarrow\mathbb Q_p\longrightarrow0.
 \end{gathered}
                                                        \tag{FF.41.5}
\]
The extension class of the last
sequence is the local Kummer class
\(\kappa_M(q)\). This use of \(q\)
is entirely local: we do not assert
that \(q\) is algebraic over
\(\mathbb Q\).

Every element of \(I'\) lies in
\(G_{M_0}\), since the quadratic
character factors through an abelian
quotient. Its cyclotomic character
is one. In a basis adapted to
(FF.41.5), the image of \(I'\)
therefore consists of unipotent
upper-triangular matrices. They
commute, so \(H=I''\) acts trivially
on \(V_pA\). It also acts trivially
on \(\mathbb Q_p(1)\). Thus (FF.41.2)
is the full two-dimensional Hom space.
Dualizing and twisting (FF.41.5) gives
\[
 0\longrightarrow\mathbb Q_p(1)
  \longrightarrow S
  \xrightarrow{\ r\ }\mathbb Q_p
  \longrightarrow0,
                                                        \tag{FF.41.6}
\]
where \(r\) is restriction to the
torus Tate line. Its extension class
is \(\kappa_M(q)\), up to sign.

For any finite \(M/M_0\), valuation
shows that \(\kappa_M(q)\) is nonzero
in \(H^1(G_M,\mathbb Q_p(1))\):
the valuation component in the
rationalized Kummer completion is
\(v_M(q)\ne0\). Restriction to
\(I_M\) is injective, by
inflation--restriction and
\(\mathbb Q_p(1)^{I_M}=0\).
Consequently (FF.41.6) does not split
on \(I_M\). Its connecting map
from \(\mathbb Q_p\) is nonzero,
and the left term has no inertia
invariants. Therefore
\[
                         S^{I_M}=0.
                                                        \tag{FF.41.7}
\]

We now determine the invisible space,
with all finite extensions included.
If \(r(\lambda)=0\), take any
finite \(E/B_p\) and \(a\in A(E)\).
Choose a number field \(D\subset E\)
containing \(D_0\) and the coordinates
of \(a\), and put \(M=D_v\).
Then \(G_E\subset G_M\), using
the chosen henselian valuation.
By (FF.41.5), write
\(a=\varphi(u)\), \(u\in M^\times\).
Choose roots \(u_n\) recursively
with \(u_{n+1}^p=u_n\) and
\(u_n^{p^n}=u\). The points
\(\varphi(u_n)\) are coherent
\(p\)-power divisions of \(a\).
For \(g\in G_M\), their Kummer
cocycle is
\[
       g\varphi(u_n)-\varphi(u_n)
           =\varphi\bigl(g(u_n)/u_n\bigr),
       \qquad g(u_n)/u_n\in\mu_{p^n}.
                                                        \tag{FF.41.8}
\]
It lies in the torus subgroup at
every level, compatibly. Hence
\(\lambda\) kills the full Tate-module
Kummer class, not merely its separate
finite-level images. Restricting to
\(G_E\) proves
\(\ker(r)\subset N\).

Conversely let \(r(\lambda)\ne0\).
We construct an actual point of
\(A(B_p)\) which detects it.
The group \(H\) is itself a closed
subgroup of wild inertia, because
\(I/P\) is abelian. It is nontrivial
by Claim~\ref{clm:FFmetabeliantest}.
Closed subgroups of a free pro-\(p\)
group are free pro-\(p\), so \(H\)
has a nontrivial continuous quotient
\(\mathbb F_p\). This uses freeness
of \(H\) itself, not merely the
presence of a free subgroup.
Kummer theory and \(\mu_p\subset B_p\)
give \(a\in B_p^\times\) which is
not a \(p\)-th power.
Since \(\mu_{p^\infty}\subset B_p\),
its compatible Kummer class is a
homomorphism
\[
       \kappa_H(a):H\longrightarrow\mathbb Z_p(1)
                                                        \tag{FF.41.9}
\]
whose reduction modulo \(p\) is
nonzero. It remains nonzero after
rationalization.

We may normalize \(a\) to an
arbitrarily small principal unit
without losing this property.
All compatible roots of \(p\)
lie in \(B_p\): prime-to-\(p\)
roots lie in its tame subfield,
and the radical tower over
\(K_p^{\rm ab}\) is metabelian
over \(K_p\), as in (FF.40.3).
For any \(s\in\mathbb Q\), a
chosen \(p^s\) has zero rational
Kummer class on \(H\), since an
integer multiple of that class
is a multiple of \(\kappa_H(p)=0\).
Divide \(a\) by \(p^{v_p(a)}\),
and divide its residue by the
matching prime-to-\(p\) root of
unity in \(K_p\). The result
\(u_0\in B_p\) has positive
valuation of \(u_0-1\) and the
same nonzero rational Kummer class.
Raising it to a sufficiently large
\(p\)-power produces
\(u=u_0^{p^m}\) as close to one
as desired, still with nonzero
rational Kummer class.

Choose a number field
\(D\subset B_p\) containing \(D_0\)
and \(u\), and put \(M=D_v\).
The auxiliary analytic point
\(a_0=\varphi(u)\in A(M)\) lies
arbitrarily close to the identity.
Equations (FF.41.8)--(FF.41.9) give
\[
        \lambda_*\kappa_H(a_0)
              =r(\lambda)\kappa_H(u)\ne0.
                                                        \tag{FF.41.10}
\]
Here \(a_0\) need not have globally
algebraic coordinates, and is only
an intermediate local point.

After multiplying \(\lambda\) by
a nonzero integer, it is integral
on \(T_pA\). Choose \(n\) such
that the integral class in
(FF.41.10) has nonzero image modulo
\(p^n\). The subgroup
\([p^n]A(M)\) is open: in a
characteristic-zero formal coordinate
at the identity, multiplication by
\(p^n\) has nonzero derivative and
is an analytic isomorphism between
sufficiently small neighborhoods.
Thus the finite-level Kummer class
is constant on an open neighborhood
of \(a_0\).

An etale coordinate at the rational
identity, defined over \(D\), shows
that \(A(D^h)\) is dense in the
small identity neighborhood of
\(A(M)\). Indeed choose its coordinate
in the dense subfield \(D\); Hensel's
lemma supplies the inverse point
over the chosen henselization
\(D^h\). Since \(B_p\) is henselian
and contains \(D\), this chosen
\(D^h\) lies in \(B_p\). We can
therefore choose
\(a'\in A(D^h)\subset A(B_p)\)
with \(a'-a_0\in[p^n]A(M)\).
The nonzero finite-level evaluation
in (FF.41.10) persists for \(a'\).
On \(H\), all the representations
are trivial, so a nonzero integral
Kummer homomorphism into the
torsion-free module \(\mathbb Z_p(1)\)
remains nonzero rationally.
Thus \(\lambda\notin N\).

We have proved \(N=\ker(r)\).
Equations (FF.41.6)--(FF.41.7)
now give all of (FF.41.3).
The accompanying standard-library program
checks the invariant polynomial identity
(FF.41.4) and the three unitriangular
commutator identities of
Claim~\ref{clm:FFmetabeliantest} over
\(\mathbb Z\), hence at every prime.
It does not compute a global Tate
parameter or certify the infinite arguments.
\end{proof}

The most direct ghost attempt from
this example would assign to a
point its multiplicative Tate
uniformization parameter. Two
requirements remain unproved.
First the analytic parameter of
a globally algebraic point need
not be known to lie in the global
field \(B_p\); the uniformization
theorem only places it in a finite
\(p\)-adic completion. Second, a
homomorphism to the geometric
Tate module is not a map of full
etale topoi. Such a map must also
give compatible maps after removing
every finite set of target points.
A chosen analytic fundamental
domain does not supply these
Zariski-open restrictions. We have
constructed neither the required
actual point function nor this
compatible family of morphisms.
The constant-point construction is
already excluded by
Claim~\ref{clm:FFstrictdetection}.
